% Part: incompleteness
% Chapter: representability-in-q
% Section: comp-representable

\documentclass[../../../include/open-logic-section]{subfiles}

\begin{document}

\olfileid{inc}{req}{crq}
\olsection{સંગણનીય વિધેયો~$\Th{Q}$માં નિરૂપ્ય છે}

\begin{thm}
દરેક સંગણનીય વિધેય
$\Th{Q}$માં નિરૂપ્ય છે.
\end{thm}

\begin{proof}
ચોક્કસતા માટે, ચર્ચ--ટ્યુરિંગ
માન્યતાનો ઉપયોગ કરીને એવું
કહીએ કે વિધેય સંગણનીય છે
ત્યારે અને માત્ર ત્યારે જ
તે સામાન્ય પુનરાવર્તી છે.
સામાન્ય પુનરાવર્તી વિધેયો
એવાં છે જે શૂન્ય વિધેય~$\Zero$,
ઉત્તરગામી વિધેય~$\Succ$ અને
પ્રક્ષેપણ વિધેય~$\Proj{n}{i}$માંથી
સંયોજન, આદિમ પુનરાવર્તન અને
નિયમિત લઘુત્તમીકરણ વડે
વ્યાખ્યાયિત થઈ શકે છે.
\olref[pri]{lem:prim-rec} પ્રમાણે,
$f$ અને~$g$માંથી વ્યાખ્યાયિત
કરી શકાય એવું કોઈપણ વિધેય~$h$,
$f$, $g$ અને $\Zero$,
$\Succ$, $\Proj{n}{i}$,
$\Add$, $\Mult$, $\Char{=}$માંથી
સંયોજન અને નિયમિત લઘુત્તમીકરણ
વડે પણ વ્યાખ્યાયિત થઈ શકે છે.
પરિણામે, વિધેય સામાન્ય પુનરાવર્તી
હોય ત્યારે અને માત્ર ત્યારે જ
તે $\Zero$, $\Succ$,
$\Proj{n}{i}$, $\Add$,
$\Mult$, $\Char{=}$માંથી
સંયોજન અને નિયમિત લઘુત્તમીકરણ
વડે વ્યાખ્યાયિત થઈ શકે છે.

વધુમાં, સંબંધિત મૂળભૂત વિધેયો
$\Th{Q}$માં નિરૂપ્ય છે તે આપણે
બતાવ્યું છે
(\cref{inc:req:bre:prop:rep-zero,inc:req:bre:prop:rep-succ,inc:req:bre:prop:rep-proj,inc:req:bre:prop:rep-id,inc:req:bre:prop:rep-add,inc:req:bre:prop:rep-mult}).
નિરૂપ્ય વિધેયોમાંથી સંયોજન
અથવા નિયમિત લઘુત્તમીકરણ
વડે વ્યાખ્યાયિત કોઈપણ વિધેય
પણ નિરૂપ્ય છે
(\olref[cmp]{prop:rep-composition},
\olref[min]{prop:rep-minimization}).
તેથી દરેક સામાન્ય પુનરાવર્તી
વિધેય~$\Th{Q}$માં નિરૂપ્ય છે.
\end{proof}

\begin{explain}
આપણે બતાવ્યું છે કે સંગણનીય
વિધેયોનો ગણ એટલે $\Th{Q}$માં
નિરૂપ્ય વિધેયોનો ગણ.
હકીકતમાં આ સાબિતી વધુ સામાન્ય
છે. નિરૂપ્યતાની વ્યાખ્યામાંથી
જોવું મુશ્કેલ નથી કે
$\Th{Q}$નું વિસ્તરણ કરતો
કોઈપણ સિદ્ધાંત (અથવા જેમાં
$\Th{Q}$નું અર્થઘટન કરી શકાય)
સંગણનીય વિધેયોને નિરૂપી શકે છે.
પણ બીજી દિશામાં, જે
!!{derivation} પદ્ધતિમાં
!!{derivation}ની વિભાવના સંગણનીય
છે, તેમાં દરેક નિરૂપ્ય વિધેય
સંગણનીય છે. આથી, દાખલા તરીકે,
સંગણનીય વિધેયોનો ગણ પીઆનો
અંકગણિતમાં, અથવા ઝર્મેલો--ફ્રેન્કેલ
સમૂહસિદ્ધાંતમાં પણ નિરૂપ્ય
વિધેયોના ગણ તરીકે વર્ણવી શકાય
છે. ગડેલે નોંધ્યું તેમ, આ કંઈક
આશ્ચર્યજનક છે. આગળ જોઈશું
કે સાબિતિક્ષમતાના પ્રશ્નો માટે
કયો સિદ્ધાંત લેવાય છે તે ખૂબ
મહત્ત્વનું છે: સ્થૂળ રીતે કહીએ
તો સ્વયંસિદ્ધિઓ જેટલી વધુ
શક્તિશાળી, તેટલું વધુ સાબિત
કરી શકાય. પરંતુ અનેક પ્રકારના
સ્વયંસિદ્ધિમૂલક સિદ્ધાંતોમાં
નિરૂપ્ય વિધેયો બરાબર સંગણનીય
વિધેયો જ છે; સિદ્ધાંતોને
અસરકારક રીતે સ્વયંસિદ્ધિઓથી વર્ણવી શકાય
ત્યાં સુધી વધુ શક્તિશાળી
સિદ્ધાંત વધુ વિધેયોને નિરૂપતો
નથી.
\sourcecorrection{OLINC-031}{સ્થિર અંગ્રેજી
મૂળના અંતિમ વાક્યમાં માત્ર
``સ્વયંસિદ્ધીકરણ શક્ય'' હોવાની
શરત લખાઈ છે. અગાઉની દલીલ
સાબિતીઓ તપાસી શકાય એવી
નિગમન-પદ્ધતિ પર આધાર રાખે છે;
તેથી અહીં અસરકારક સ્વયંસિદ્ધીકરણની
શરત સ્પષ્ટ કરી છે.}
\end{explain}

\end{document}
